% Separate post-CP23 candidate. No conclusion follows from finite row checks alone.
\subsection{A uniform middle-ramification quotient argument}

Fix $p\ge5$ and put
\[
 N=p^2-1,\qquad e_0=N/2,\qquad s=e_0(p+1).
                                                        \tag{B-RD66.1}
\]
After completed maximal unramified base change, $K$ is complete of
mixed characteristic $(0,p)$ with algebraically closed residue field,
$e_K=e_0$ and $\mu_p\subset K$. For a pair of whole potentially
supersingular height-two groups take a common tame good field
$L/K$ of degree $N$. Thus $v_L(p)=e_0N=(p-1)s$.
In the nonsquare finite-determinant case normalize by a common
finite height-one twist as in B50, so an actual special tangent
exponent $c$ satisfies $c\equiv1\pmod{p-1}$. Choose
$\sigma\pi_L=\zeta\pi_L$, with $\zeta$ primitive of order $N$.

\begin{lemma}[The unequal Newton comparison rules at middle ramification]
\label{lem:B-RD66-rules}
\status{Candidate; full general-index Newton and point-field proof written.}
A reducible normalized residual module has a unique high line $H$
and an integer $B$ with
\[
 0<B<pe_0,\quad B\equiv c\pmod{p+1},\quad
 \alpha=s-B,\quad\beta=B/p,\quad\alpha>\beta.          \tag{B-RD66.2}
\]
Here nonzero points of $H$ have valuation $\alpha$, all other
nonzero $p$-torsion points have valuation $\beta$, and the high-line
character is $\zeta^{c-B}$. The integer $B$ is odd. If $P$ has
order $p^2$ and $pP\notin H$, then $v_L(T(P))=B/p^3$.
For two such groups with isomorphic actual $p^2$-torsion:
\begin{enumerate}
\item In the residual wild case their high lines agree. If both $B_i$
are prime to $p$, they are equal and the special characters agree.
Exactly one $B_i$ cannot be divisible by $p$.
\item In the residual tame case write $B_i=pb_i$. If both $b_i$ are
prime to $p$, then $b_1=b_2$, their high lines agree, and the special
characters agree.
\end{enumerate}
\end{lemma}
\begin{proof}
The three Newton vertices of $[p](T)$ are
$(1,(p-1)s),(p,b),(p^2,0)$ with $b=v_L(u)$, using exactly the
universal coefficient bounds proved in Lemma~\ref{lem:B-RD49-point}.
If the polygon has one slope, its leading-coordinate residual
representation is tame and either irreducible over $\mathbf F_p$
or scalar. The scalar case has square determinant, so is excluded.
Reducibility therefore gives two slopes. The high line is inertia
stable and its character takes values in $\mathbf F_p^\times$;
it is rational over $L$. Consequently its valuation is integral,
so $b=(p-1)B$ with $B\in\mathbf Z$. The two slopes give
$\alpha=s-B$ and $\beta=B/p$, and their strict order gives
$B<ps/(p+1)=pe_0$. Exact tame averaging gives
$b\equiv(p-1)c\pmod N$, hence $B\equiv c\pmod{p+1}$.
Since $s=((p+1)/2)N$ is a multiple of $N$, the high-line character
is $\zeta^{c+s-B}=\zeta^{c-B}$. In particular $B$ is odd.
For target valuation $\beta$, the inverse Newton polygon has one
segment of degree $p^2$, giving $B/p^3$ for every preimage.

For completeness the comparison rules are reproved at this index.
The common residual division field $R/L$ is trivial or cyclic of
order $p$. In the wild case its high line is its unique fixed line.
For a low point, $v_R(T(Q))=B$ and its difference under a nontrivial
wild element has valuation $p\alpha=p(s-B)$. The exact lower-break
lemma~\ref{lem:B-RD60-break} gives
\[
 r=ps-(p+1)B\ (p\nmid B),\qquad
 r<ps-(p+1)B\ (p\mid B).                              \tag{B-RD66.3}
\]
Two prime-to-$p$ parameters are equal from the same actual break;
their high-line characters then identify $c$ modulo $N$.
If only $B_2$ is divisible by $p$, the common break forces
$B_1>B_2$. Choose corresponding $P_i$ with $pP_i$ outside the
common high line. Over $R$, their valuations are $B_i/p^2$.
The prime-to-$p$ denominator on the first side and the elementary
point-field lemma force the common $E=R(P_i)$ to have degree $p^2$;
all translations in $M[p]$ occur. Thus $[E:L]=p^3$ and the point
coordinates have valuations $B_i$. A translation outside the high
line has difference valuations $p^2B_i$. Its actual break from the
first coordinate is $NB_1$, while the second coordinate would require
$p^2B_2\ge B_2+NB_1$, a contradiction.

In the tame case $R=L$, so $B_i=pb_i$. Choose $pP_i$ outside both
high lines. If both $b_i$ are prime to $p$, their common point field
has degree $p^2$, and their coordinate valuations are $b_i$.
A translation outside both high lines gives the same break
$Nb_1=Nb_2$. Thus $b_1=b_2=b$. A translation high only in one
model, if it existed, would force the strictly higher valuation
$\alpha$ to equal $\beta=b$ by the same prime-valuation equality.
The high lines therefore agree and their characters give $c_1=c_2$.
This proves both rules without a minimal-ramification assumption.
\end{proof}

\begin{lemma}[High and low cyclic quotient coordinates]
\label{lem:B-RD66-quotients}
\status{Candidate; complete integral-coordinate and torsion-quotient proof written.}
For a group as in Lemma~\ref{lem:B-RD66-rules}, quotient by its
high line transforms its Newton parameter by
\[
 B_{\rm high}=\begin{cases}pB,&B<e_0,\\s-B,&B>e_0.\end{cases}
                                                        \tag{B-RD66.4}
\]
The equality case is impossible since $e_0$ is even and $B$ odd.
If the full residual module is rational over $L$ and $J$ is an
inertia-stable low line, quotient by $J$ transforms it by
\[
                         B_{\rm low}=B/p.             \tag{B-RD66.5}
\]
These are actual good formal quotients over $\mathcal O_L$;
their generic quotients descend to $K$. Their special tangent
exponent is $pc$ modulo $N$, so their finite special character
is unchanged. An actual $p^3$-torsion isomorphism preserving the
chosen lines induces an actual $p^2$-torsion isomorphism on the
quotients.
\end{lemma}
\begin{proof}
The finite flat closure of a chosen line has rank $p$. The good
formal quotient exists, is isogenous to the original height-two
group, and its parameter map is a unit times
$\prod_{Q\in J}(T-T(Q))$ by Weierstrass preparation. Thus the
quotient coordinate valuation is the sum of these difference
valuations. All statements use integral quotient coordinates.

For $J=H$, a point of valuation less than $\alpha$ acquires
valuation multiplied by $p$. The line $M[p]/H$ therefore has
valuation $B$. The other points arise from preimages of nonzero
high points. If $B>e_0$, the inverse Newton polygon has one segment,
with all roots of valuation $\alpha/p^2$, giving $p^2-p$ quotient
points of valuation $\alpha/p$. Since $B>\alpha/p$, this is
$B_{\rm high}=s-B$.
If $B<e_0$, the inverse polygon has $p$ roots of valuation
$(s-pB)/p$ and $p^2-p$ roots of valuation $B/p$.
Both are less than $\alpha$. Taking all $p-1$ nonzero high targets
and dividing the preimage counts by $p$ gives $p-1$ quotient points
of valuation $s-pB$ and $p^2-p$ of valuation $B$.
This gives $B_{\rm high}=pB$.

For a low line $J$, all its nonzero points have valuation $\beta$.
Every nonzero class in $M[p]/J$ has a unique high representative,
so its quotient valuation is
$\alpha+(p-1)\beta=s-B/p$. The other quotient $p$-torsion points
arise from preimages of nonzero low points. They have source
valuation $\beta/p^2<\beta$, hence quotient valuation
$\beta/p=B/p^2$. Their number is $p^2-p$, yielding (B-RD66.5).
Here $B/p$ is an integer because $M[p]$ is rational over $L$.

The special rank-$p$ subgroup of a one-dimensional connected
formal group is its Frobenius kernel. The special quotient is
therefore the Frobenius twist with inherited tangent action $pc$,
and its rational special character is unchanged. The generic
quotient descends because the chosen line was inertia stable.
Finally the inverse image of the quotient's $p^2$-torsion is
$\{P\in M[p^3]:p^2P\in J\}$; quotienting this actual finite group
by $J$ proves the asserted induced torsion isomorphism.
\end{proof}

\begin{proposition}[Nonsquare formal recovery at middle ramification]
\label{prop:B-RD66-nonsquare}
\status{Candidate; complete case and quotient proof written.}
Two normalized whole potentially supersingular height-two groups
with isomorphic actual $p^3$-torsion have the same finite special
inertia character.
\end{proposition}
\begin{proof}
For an irreducible residual module both Newton polygons have one
slope. The leading character is obtained from the special tangent
exponent by adding $e_0=N/2$, a scalar sign in $\mathbf F_p^\times$.
Since $pe_0\equiv e_0\pmod N$, the residual Frobenius orbit
$\{c+e_0,p(c+e_0)\}$ determines $\{c,pc\}$.
A reducible one-slope module has square determinant, so the remaining
cases have the unequal parameters of Lemma~\ref{lem:B-RD66-rules}.

If residual wild inertia is nontrivial, the high line is common.
The rules already handle every case where at least one $B_i$ is
prime to $p$. Suppose both are divisible. If both $B_i<e_0$, put
$B_i=pb_i$. Then $0<b_i<e_0/p<p/2$. The common high-line character
and $c_i\equiv1\pmod{p-1}$ imply
$b_1\equiv b_2\pmod{p-1}$, hence $b_1=b_2$ and $c_1=c_2$ modulo $N$.
If both $B_i>e_0$, quotient by their common high line. The parameters
$s-B_i$ are prime to $p$ because $p\nmid s$, and the actual
$p^2$-torsion remains isomorphic. The rules make the quotient
parameters and characters equal. If one $B_i$ lies on each side
of $e_0$, the quotient parameters have opposite $p$-divisibility,
so the rules forbid their $p^2$-torsion isomorphism. This exhausts
the wild case.

Now let residual inertia be tame, and write $B_i=pb_i$ with
$0<b_i<e_0$. If both $b_i$ are prime to $p$, the rules identify the
characters. Suppose first that $b_1=b$ is prime to $p$ and
$b_2=pk$. Then $0<k<e_0/p<p/2$ and $B_2=p^2k>e_0$.
Choose corresponding order-$p^2$ points with residual targets
outside both high lines. Their common point field has degree $p^2$
over $L$, forced by the coordinate valuation $b/p^2$ in the first
model. All translations occur, and the normalized point valuations
are $b$ and $pk$.
If their high lines differed, choose a translation high in the first
model and low in the second. The prime-valuation coordinate gives
actual break $p^2\alpha_1-b$, while the lower-bound inequality for
the other would require
\[
                  p^2(\alpha_1-pk)\le b-pk.            \tag{B-RD66.6}
\]
But $\alpha_1=s-pb\ge e_0+p$, $pk<e_0$ and $b<e_0$;
the left side is greater than $p^3$, while the right side is less
than $e_0<p^2$. Thus the high lines coincide.
Quotient by that common line. The second parameter $s-p^2k$ is
prime to $p$. If $B_1<e_0$, the first quotient parameter $pB_1$
is divisible by $p$, impossible by the rules and the inherited
$p^2$-torsion isomorphism. If $B_1>e_0$, both new parameters are
prime to $p$; their equality forces $B_1=B_2$, contrary to
$p\nmid b$ and $b_2=pk$. Thus the mixed $b$-divisibility case is
impossible.

Finally suppose $b_i=pk_i$ for both groups. Then
$B_i=p^2k_i>e_0$. If their high lines agree, quotienting by that
line gives two prime-to-$p$ parameters $s-B_i$, which the rules
identify. If their high lines differ, quotient the first group
by its high line and the second by the corresponding low line.
Both chosen lines are inertia stable. The quotient parameters are
$s-B_1$, prime to $p$, and $B_2/p=pk_2$, divisible by $p$.
Their actual $p^2$-torsion isomorphism is forbidden by the rules.
These alternatives exhaust tame residual inertia. Isogeny invariance
in every quotient comparison proves the original character equality.
\end{proof}

\begin{theorem}[A uniform additional exact least-index family]
\label{thm:B-RD66}
\status{Candidate; full square, filtered and mixed reductions written.}
For every $p\ge5$, in the full B.4 monodromy setting with height two,
perfect residue field and $\mu_p\subset K$,
\[
                    \nu_{\min}\bigl(p,(p^2-1)/2,2\bigr)=2.
                                                        \tag{B-RD66.7}
\]
\end{theorem}
\begin{proof}
The lower example is Theorem~\ref{thm:B-RD63-arc} with $n=2$,
since $e_0>p$ and $p-1$ divides $e_0$.
For sufficiency pass to the completed maximal unramified base.
Since $p\nmid e_0$, the cyclotomic image modulo $p^3$ has order $p^2$.
The determinant and dimension argument of Theorem~\ref{thm:B-RD64}
reduces to two whole potentially supersingular groups with a common
finite determinant: dimension zero/two and two filtered systems are
already determined; the mixed comparison is impossible by
Proposition~\ref{prop:B-RD63-mixed}, since $e_0<p^3$.
The nonsquare finite determinant case is the preceding proposition.

If the finite determinant is square, twist both groups by its finite
square root. Their common good field has tame degree $m\mid p+1$,
so $a=v_L(p)\le e_0(p+1)$. Different special characters again give
an integral tame projector and the self-orthogonal graph closure
$\ker F_M^3\times\ker F_N^3$, exactly as in
Lemma~\ref{lem:B-RD65-square}. The actual target $Q\times M$ of
the no-flat-kernel map there has product height-two special fibre.
Two applications of $[p]$ to a coefficientwise valuation-one vector
have valuations at least
\[
 q_1\ge p+1,\qquad q_2\ge\min\{a+p+1,p^2+p+1\}\ge a/(p-1).
                                                        \tag{B-RD66.8}
\]
Indeed $a\ge e_0\ge p$, and
\[
     a/(p-1)\le(p+1)^2/2\le p^2+p+1.
\]
The closed logarithm and actual Cartier saturation proof already
established in B64 forces $F^2$ to annihilate the special kernel
$\ker F_N^3$, which is impossible. This is an actual-target bound,
not inherited Hasse thickness. Thus square and nonsquare cases
both recover the character at actual $p^3$-torsion. Together with
the lower example this proves the exact manuscript index.
\end{proof}

\begin{remark}[Scope of the uniform quotient argument]
\label{rem:B-RD66-scope}
\status{Candidate scope; unrelated indices remain open.}
The family includes B65 at $p=5,e_0=12$, and for example gives
$e_0=24$ at $p=7$. It does not settle $p=5,e=16$. The proof uses
both high and low stable cyclic quotients, with their actual
$p^3$-to-$p^2$ comparison, and a strict bound on the small integer
$B/p$; no enumeration in finitely many primes proves the theorem.
\end{remark}
